%==============================================================================
\section{Box and Whisker Plots --- Biểu đồ hộp và râu}
\label{sec:box-whisker}
%==============================================================================

\begin{dinhnghia}[title={Boxplot (Box and Whisker Plot)}]
Boxplot là biểu diễn đồ họa tóm tắt dữ liệu bằng \textbf{năm số thống kê}:
giá trị nhỏ nhất, tứ phân vị thứ nhất (Q1), median, tứ phân vị thứ ba (Q3), giá trị lớn nhất (trong phạm vi ``râu'' chuẩn).
\end{dinhnghia}

\begin{ynghia}[title={Cấu trúc boxplot}]
\begin{itemize}
  \item \textbf{Hộp (box)}: khoảng liên tứ phân vị (IQR) từ Q1 đến Q3.
  \item \textbf{Râu (whiskers)}: kéo từ hộp tới giá trị cao/thấp nhất còn nằm trong $1{,}5 \times \mathrm{IQR}$ tính từ biên hộp.
  \item \textbf{Điểm ngoài râu}: outlier.
\end{itemize}
\end{ynghia}

\subsection{Boxplot trong Python}

\begin{phuongphap}[title={Dataset Iris}]
Ví dụ dùng \textbf{Iris} (scikit-learn): 150 hoa, 4 đo (sepal length/width, petal length/width), 3 loài Setosa, Versicolor, Virginica.
\end{phuongphap}

\subsection{Ví dụ --- Boxplot theo loài}

\begin{vidu}[title={Nạp dữ liệu}]
\begin{lstlisting}
import matplotlib.pyplot as plt
import seaborn as sns
from sklearn.datasets import load_iris
iris = load_iris()
data = iris.data
target = iris.target
\end{lstlisting}
\end{vidu}

\begin{vidu}[title={Boxplot sepal length theo species}]
\begin{lstlisting}
import matplotlib.pyplot as plt
import seaborn as sns
from sklearn.datasets import load_iris
iris = load_iris()
data = iris.data
target = iris.target

plt.figure(figsize=(7.5, 3.5))
sns.boxplot(x=target, y=data[:, 0])
plt.xlabel('Species')
plt.ylabel('Sepal Length (cm)')
plt.show()
\end{lstlisting}

Trục $x$: chỉ số loài (0, 1, 2); trục $y$: sepal length (cm).
\end{vidu}

\begin{ynghia}[title={Đọc biểu đồ}]
Setosa có sepal length \textbf{ngắn hơn} Versicolor và Virginica (median và IQR thấp hơn).
Versicolor và Virginica có median và độ trải tương tự.
Setosa hầu như không có outlier; Versicolor và Virginica có vài điểm ngoài râu.
\end{ynghia}

\begin{vidu}[title={Output}]
\tsimg{04_box_whisker_plots/img01.jpg}
\end{vidu}
